\section*{VI. Kết luận}

Nghiên cứu đã xây dựng mô hình toán học giải tích dạng đóng cho mạng SAGIN-NOMA hai chặng trên kênh fading Rician tổng quát thông qua hàm Marcum-Q bậc 1, được kiểm chứng chính xác bằng mô phỏng Monte-Carlo $10^5$ mẫu. Khung tối ưu hóa luân phiên (AO) đã thiết lập nghiệm tham chiếu số học chuẩn, làm giảm xác suất dừng toàn phần từ $0.0401$ xuống $0.0206$ (giảm $48.6\%$ so với quỹ đạo thẳng, và giảm $57.2\%$ so với kịch bản không tối ưu) trên hành trình $N=10$ khe thời gian. Đặc biệt, tác tử Học tăng cường sâu Actor-Critic (DDPG) đã chứng minh khả năng đạt hiệu năng xác suất dừng tương đương ($0.0197$) trong khi rút ngắn thời gian tính toán ra quyết định xuống chỉ còn $\mathbf{1.5}$ \textbf{ms/khe} (nhanh hơn $120$ lần so với AO), đồng thời bảo đảm an toàn ranh giới SIC $100\%$ với vùng đệm tự nhiên cách ranh giới $0.06 - 0.10$. Kết quả này khẳng định tính khả thi và tiềm năng ứng dụng to lớn của mạng học sâu nhúng trong điều khiển bay tự hành trực tuyến cho các hệ thống thông tin không dây thế hệ mới 6G.
